\scanpage{3}
\uline{Catégorie des motifs sur $\mathbf F_p=k$}

$\mathcal M$ admet un foncteur fibre $T_\ell$ sur $Q_\ell$ (à valeurs en dans $\operatorname{Mod}(\widehat{\mathbf Z},Q_\ell)$) pour tt $\ell\ne p$.
Pour $\ell=p$, il y a un foncteur $T_p$ sur $Q_p$, cohomologie cristalline (foncteur à valeurs dans \deleted{$\operatorname{Mod}($}$F$-cristaux sur $Q_p$, donc dans vectoriels sur $Q_p$ munis d'un autom., donc dans $\operatorname{Mod}(\mathbf Z,Q_p)$ --- Attention, ici on ne peut remplacer $\mathbf Z$ par $\widehat{\mathbf Z}$ !]

Soit \deleted{$\unsure{\mathcal M}$} $\operatorname{Niv}^0(\mathcal M)$ la sous-catégorie des motifs de niveau $0$
\[
\operatorname{Niv}^0(\mathcal M)^0\qquad\rule{6cm}{0.3pt}\quad\rule{0.7cm}{0.3pt},\ \text{poids zéro.}
\]
Il y a un foncteur fibre naturel \deleted{\unsure{indépendant}} \inserted{intrinsèque} en $Q$ sur la sous-catégorie $\operatorname{Niv}^0(\mathcal M)^0$, soit
\diagram{\begin{tikzcd}
F^0:\operatorname{Niv}^0(\mathcal M)^0 \arrow[rr] \arrow[dr] && \operatorname{Mod}(Q,\deleted{\widehat{\mathbf Z}}) \\
& \operatorname{Mod}(Q,\widehat{\mathbf Z}) \arrow[ur] &
\end{tikzcd}}
\margin{\unsure{équivalent} aux faisceaux de $Q$-(\unsure{vect.}\ \unsure{dim. finies}) sur $k$}
\margin{\unsure{valable} pour tout $k$}
défini par
\[
F^0(H^*(X))=f_*(\uline{Q}_X)
\]
si $X$ fini sur \deleted{$k$} --- cela définit $F^0$ sur $\operatorname{Niv}^0(\mathcal M)^0$
\[
F^0(Q(1))=Q.
\]
Donc la gerbe \inserted{$\mathcal G$} des foncteurs fibres de $\mathcal M$ contient la sous-gerbe \inserted{\deleted{$\unsure{\mathcal G}$}$\unsure{\mathcal H}$} des foncteurs fibres de $\mathcal M$ prolongeant $F^0$, laquelle correspond (si $G$ est le lien de $\mathcal M$, qui est \inserted{\unsure{un}} un groupe \uline{commutatif} sur $k$, si on admet les \unsure{p\textsuperscript{tés}} de \unsure{densité|dualité}....) au groupe
\[
H=\operatorname{Ker}(G^0\xrightarrow{\unsure{\epsilon|e}}G_m)
\]
Par ailleurs, \deleted{\unsure{si on admet}} en vertu des conj. de Tate et de Tate-Honda, $G$ est le groupe de \unsure{t. m.} sur $Q$, dont le groupe des caractères $\Gamma$ est
\mspage{2}
